% !TeX root = ../thuthesis-example.tex
\chapter{系统设计与实现}

\section{系统总体架构设计}

\subsection{系统模块组成}
从当前项目实现来看，系统启动后先进入登录窗口，登录成功后进入主窗口。主窗口以内嵌标签页的方式组织多个页面，包括番茄钟页面、正向计时页面、数据统计页面、记账页面、日记页面，以及图片处理、设置与同步等扩展页面。这样的组织方式既保证了主业务功能的集中展示，也为后续系统扩展预留了足够空间。
在窗口组织方式上，主窗口采用QMainWindow作为顶层容器，内部通过QTabWidget实现多页面切换。与为每项功能单独创建窗口相比，这种方案能够显著降低用户在多个界面之间频繁切换的成本，也使整体交互更加集中。与此同时，工具栏中提供日间主题与夜间主题切换动作，进一步提升了系统的一体化体验。对毕业设计项目而言，这种窗口组织方式结构清晰、效果直观，便于展示整体系统架构。

\subsection{页面关系与业务流程}
用户登录后，会话模块保存当前用户名，后续各页面均通过统一会话对象获取当前用户信息。番茄计时和正向计时模块在完成记录后，会将专注信息写入用户目录下的日期文本文件中；日记模块负责同一日期文本内容的读取与维护；记账模块则将账单对象持久化到独立的账单文件中；统计模块通过读取专注记录和账单记录，对周、月等范围内的数据进行汇总并可视化展示。由此形成“用户登录—业务记录—数据落盘—统计分析—图形展示”的完整业务闭环。
进一步来看，该业务流程体现出较为典型的前后分层结构。界面页面负责收集用户输入和呈现结果，不直接处理复杂的持久化细节；服务层则集中管理路径生成、文件读写、对象序列化和统计汇总等逻辑。以记账流程为例，用户在页面中完成输入后，页面只负责校验和触发提交动作，真正的数据保存由记账服务类负责执行；统计页面也不直接维护原始数据，而是在刷新时从服务层重新加载并汇总结果。这种模式有助于控制复杂度，提高代码复用性。

\section{核心功能模块实现}

\subsection{番茄计时与正向计时模块}
番茄计时模块以QTimer为核心，通过剩余秒数、轮次状态和休息状态等变量控制计时流程，支持开始、暂停、继续和提前结束等操作。系统允许用户自定义任务名称、番茄轮数、专注时长和休息时长，并在每轮专注完成后自动写入专注记录。正向计时模块则适用于不固定时长的专注场景，系统在计时结束后将秒级数据转换为分钟并写入对应日期文件，以便后续统计模块统一分析。
在实现细节上，番茄计时模块会在每次计时推进时更新界面显示，并根据是否处于休息状态决定下一阶段的流程。当用户选择提前结束时，系统能够根据当前已完成的专注秒数进行向上取整，计算对应分钟值并记录入文件，尽量保证专注成果不会因中途中断而完全丢失。正向计时模块则更适合开放式记录场景，其逻辑相对简洁，但与番茄计时模块一样，都能够与后续统计模块形成统一的数据来源。

\subsection{日记存储与记录模块}
日记模块以日期为索引，采用“一天一文件”的策略保存用户内容。系统通过统一的数据存储工具类生成用户目录，并在该目录下按照“年-月-日.txt”的形式组织文件。该方式实现简单、结构清晰，便于用户后续进行文件级备份和迁移，也便于统计模块按日期扫描文件内容，提取专注记录和相关信息。
除了基础文本记录外，日记页面还支持字体、字号、加粗、斜体等基础编辑功能，并能够插入本地图片，从而增强记录的表达能力。页面还结合日历窗口实现日期切换，用户可以直观查看不同日期的内容。项目还预留了语音输入能力，使日记模块在功能上不仅停留在简单文本编辑，而是向更丰富的人机交互方向扩展。对毕业设计而言，这部分内容能够体现系统在“记录”功能上的完整性和可扩展性。

\subsection{记账录入与编辑删除模块}
记账模块是本文的重点实现内容之一。系统以记账页面为交互核心，提供事项、金额和类型三类输入项，并通过表格展示全部账单记录。底层采用独立的记账存储服务，将每条账单封装为包含时间、事项、金额和类型的结构体，再序列化为JSON文本逐行写入用户数据文件。为了提升可维护性，系统还对消费类型进行了独立保存，支持默认类型加载与自定义类型扩展。

在列表交互方面，系统通过表格组件显示时间、事项、金额和类型四列数据，并维护界面行号与底层记录索引的映射关系，避免在排序或刷新后出现误编辑、误删除的问题。用户右键点击记录后可弹出上下文菜单，并执行“编辑”或“删除”操作。编辑时，系统将原记录内容回填到输入区域，切换到修改状态；删除时，系统先进行确认，再刷新表格并给出提示信息。上述设计既满足了基本的CRUD需求，也增强了界面的交互友好性。
在实现本模块的过程中，系统对用户输入进行了必要校验，例如事项不能为空、金额需要可解析为数值、类型需要来自有效选项或被保存为新类型。这样的输入约束有助于减少脏数据写入，保证后续统计结果的准确性。此外，右键菜单操作的实现也体现了系统对细节体验的重视：通过QMenu提供快捷操作入口，通过QPropertyAnimation控制透明度变化，使弹出和关闭过程更加平滑。对于一款桌面工具软件而言，这类细节虽然不属于核心业务算法，但直接影响用户对系统完成度的感受。

\section{数据统计与界面交互实现}

\subsection{消费统计与图表切换}
系统数据统计模块分为专注统计和消费统计两部分。专注统计主要针对周和月两个范围，汇总番茄事件数量、正向计时事件数量以及专注总时长。消费统计则支持按日、周、月、年四种时间范围筛选账单数据，并计算总支出、平均值、最大值和最小值。为了提高展示效果，系统设计了饼状图和柱状图两类图表，并通过堆叠窗口组件实现图表视图切换，使用户可以从结构分布和数量对比两个角度观察数据结果。
专注统计和消费统计虽然来源不同，但在页面展示上采用了统一的设计思想，即“数据汇总+图形可视化+交互切换”。这种统一性有助于降低用户的理解成本，也有利于形成系统整体风格。对于消费统计模块，系统按照类型对账单金额进行聚合，并在图表中突出显示不同类别支出的占比关系；对于专注统计模块，则以时间维度展示番茄计时与正向计时的累计时长，使用户能够直观感知近期的专注投入情况。

\subsection{夜间主题与右键菜单交互}
在界面交互优化方面，系统支持日间主题和夜间主题两种视觉模式。主题切换通过统一样式管理实现，可对按钮、表格、标签及图表颜色进行整体调整，从而改善夜间环境下的阅读体验。与此同时，为了增强记账列表的操作效率，系统在右键菜单弹出过程中加入了透明度动画，使菜单出现与消失更加平滑自然，提升了界面细节质量。
从当前项目实现可以看出，主题管理并非局限于简单的背景色切换，而是同时覆盖了主窗口、对话框、工具栏、标签页、按钮、输入框、下拉框、表格、菜单以及日历控件等多类界面元素。这说明系统在视觉一致性方面进行了较细致的处理。夜间模式下，系统使用深色背景与较高对比度的文字和边框色，能够减轻用户在弱光环境中的视觉疲劳，同时避免部分控件在暗色环境下出现不可读的问题。

\subsection{本地数据持久化设计}
系统数据持久化以本地用户目录为中心展开。基础数据存储工具负责创建应用数据根目录和用户专属目录，日记与专注记录通过文本方式保存，记账记录通过JSON对象逐行保存，记账类型则单独存放为类型文件。该方案无需额外数据库环境，适合单机运行和课程设计场景，同时具有实现简单、易于调试、便于迁移的优点。
在工程实践中，这种设计还具备较好的容错性。当某一类业务数据文件损坏时，通常只会影响对应模块，而不会造成整个系统全部数据不可恢复。与此同时，按用户目录分类存放数据，也方便后续增加导出、备份和迁移等功能。虽然该方案在大规模数据处理和复杂检索方面存在不足，但就当前项目的功能定位和数据规模而言，是一种性价比较高的实现方式。